\subsection{局部凸终拓扑}

命题 5. --- 设 E 是向量空间，\((\mathrm{F}_{\alpha})_{\alpha\in\mathbf{A}}\) 是一族拓扑向量空间，且对每个 \(\alpha\in A\) ，设 \(g_{\alpha}\) 是 \(F_{\alpha}\) 到 E 中的线性映射。

\begin{enumerate}
\def\labelenumi{(\roman{enumi})}
\item
  用 \(\mathfrak{B}\) 表示满足如下条件的吸收、对称凸子集 \(\mathrm{V}\) （\(\mathrm{E}\) 中的）所成的族：\(g_{\alpha}^{-1}(\mathrm{V})\) 是 \(\mathrm{F}_{\alpha}\) 中 0 的邻域（对每个 \(\alpha\) ）；族 \(\mathfrak{B}\) 是 \(\mathrm{E}\) 上某个与向量空间结构相容的拓扑 \(\mathcal{T}\) 的 0 的基本邻域系。
\item
  线性映射 \(f\) （\(\mathbf{E}\) 到局部凸空间 \(\mathbf{G}\) 中的）（相应地，半范数 \(p\) （\(\mathbf{E}\) 上的））对 \(\mathcal{T}\) 连续，必须且只需对每个指标 \(\alpha, f \circ g_{\alpha}\) （相应地，\(p \circ g_{\alpha}\) ）在 \(\mathrm{F}_{\alpha}\) 中连续。
\item
  拓扑 \(\mathcal{T}\) 是 \(\mathbf{E}\) 上使诸 \(g_{\alpha}\) 连续的局部凸拓扑中最细的一个。
\end{enumerate}

此外，拓扑 \(\mathcal{T}\) 是 \(\mathbf{E}\) 上对线性映射（相应地，对半范数）满足条件 (ii) 的唯一的局部凸拓扑。

由于 \(\mathfrak{B}\) 是在比 \(>0\) 的位似下不变的滤子基，断言 (i) 由 II, 第 23 页, prop. 1 立即得出。由 \(\mathfrak{B}\) 的定义，拓扑 \(\mathcal{T}\) 是 E 上使诸 \(g_{\alpha}\) 连续的局部凸拓扑中最细的；由此得 (iii)。最后，显然若 \(f\) 连续，则 \(f \circ g_{\alpha}\) 也连续；反之，若诸 \(f \circ g_{\alpha}\) 对每个 \(\alpha\) 都连续，则对 G 中 0 的每个对称凸邻域 W，集 \(g_{\alpha}^{-1}(f^{-1}(\mathbf{W}))\) 是 \(\mathbf{F}_{\alpha}\) 中 0 的邻域（对每个 \(\alpha\) ）。而 \(f^{-1}(\mathbf{W})\) 是吸收、对称且凸的，因此 \(f^{-1}(\mathbf{W})\) 是 \(\mathcal{T}\) 下 0 的邻域，从而 \(f\) 连续。类似地，若 \(p\) 是 E 上的半范数，使得 \(p \circ g_{\alpha}\) 对每个 \(\alpha\) 都连续，且 U 是点 \(x \in \mathbf{E}\)（满足 \(p(x) < 1\)）所成的集，则对每个 \(\alpha\) ，集 \(g_{\alpha}^{-1}(\mathbf{U})\) 是 \(\mathbf{E}_{\alpha}\) 中 0 的对称、吸收的凸邻域；于是 U 是 E 中 0 的邻域，从而 \(p\) 连续 (II, 第 2 页, prop. 1)。

最后的陈述由 S, IV, § 2.5, 判别准则 CST 18 得出。

我们称 \(\mathcal{T}\) 为拓扑族 \(\mathcal{T}_{\alpha}\) （诸 \(F_{\alpha}\) 的）关于线性映射族 \(g_{\alpha}\) 的终局部凸拓扑。

可能发生这样的情形：\(\mathcal{T}\) 并不是 E 上与其向量空间结构相容且使诸 \(f_{\alpha}\) 连续的拓扑中最细的 (II, 第 75 页, exerc. 15；另见 II, 第 75 页, exerc. 14)。

在最重要的情形 \(E = \sum_{\alpha \in A} g_{\alpha}(F_{\alpha})\) 中，我们如下得到 T 的 0 的一个基本邻域系；对每个 \(\alpha \in A\) ，设 \(V_{\alpha}\) 是 \(T_{\alpha}\) 下 0 的一个对称邻域，作诸 \(g_{\alpha}(V_{\alpha})\) （\(\alpha \in A\)）的并，并用 \(\Gamma((g_{\alpha}(V_{\alpha})))\) 表示这个并在 E 中的凸包；由于 E 的每个元都形如 \(\sum_{\alpha \in J} x_{\alpha}\) ，其中 J 是 I 的有限子集且 \(x_{\alpha} \in g_{\alpha}(F_{\alpha})\) ，立见 \(\Gamma((g_{\alpha}(V_{\alpha})))\) 是 E 中的吸收对称凸集（每个 \(V_{\alpha}\) 在 \(F_{\alpha}\) 中都是吸收的）；由于 \(\Gamma((g_{\alpha}(V_{\alpha})))\) 包含所有 \(g_{\alpha}(V_{\alpha})\) ，它是 T 下 0 的一个邻域。另一方面，显然对 T 下 0 的每个对称凸邻域 V，有 \(V \supset \Gamma((V \cap g_{\alpha}(F_{\alpha})))\) ，我们的断言由此得证。

推论 1. --- 沿用 prop. 5 的记号，设 H 是 E 到局部凸空间 G 中的线性映射组成的一个集。假设 E 是其诸子空间 \(g_{\alpha}(F_{\alpha})\) 之和；则 H 对 T 等度连续，必须且只需对每个 \(\alpha\) ，当 f 取遍 H 时诸映射 \(f \circ g_{\alpha}\) 所成的集在 \(F_{\alpha}\) 中等度连续。

记住 I, 第 9 页, prop. 6，论证与 prop. 5 的 (ii) 类似。设 W 是 G 中 0 的一个对称凸邻域，并注意若诸映射 \(f \circ g_{\alpha}\) （其中 \(f \in \mathbf{H}\) ）所成的集等度连续，则交 \(\bigcap_{f \in \mathbf{H}} g_{\alpha}^{-1}(f^{-1}(\mathbf{W}))\) 是 \(\mathbf{F}_{\alpha}\) 中 0 的对称凸邻域。由于这个交与 \(g_{\alpha}^{-1}\big(\bigcap_{f \in \mathbf{H}} f^{-1}(\mathbf{W})\big)\) 相同，而集 \(\bigcap_{f \in \mathbf{H}} f^{-1}(\mathbf{W})\) 是对称且凸的，问题只在于证明它还是吸收的。而由假设，每个 \(x \in \mathbf{E}\) 都可写成 \(\sum_{i=1}^{n} g_{\alpha_i}(z_{\alpha_i})\) ，其中 \(z_{\alpha_i} \in \mathrm{F}_{\alpha_i}\) 。为证明存在 \(\lambda > 0\) 使 \(f(\lambda x) \in \mathbf{W}\) 对所有 \(f \in \mathbf{H}\) 成立，只须考虑 \(x = g_{\alpha}(z_{\alpha})\) （其中 \(z_{\alpha} \in \mathrm{F}_{\alpha}\) ）的情形（因为把 \(\mathbf{W}\) 换成 \(\mathbf{W}/n\) 即可过渡到一般情形）。而这种情形可由 \(g_{\alpha}^{-1}\big(\bigcap_{f \in \mathbf{H}} f^{-1}(\mathbf{W})\big)\) 是 \(\mathrm{F}_{\alpha}\) 中 0 的邻域这一事实得出。

推论 2. --- 设 \((\mathbf{J}_{\lambda})_{\lambda \in \mathbf{L}}\) 是指标集 A 的一个划分。设 \((\mathbf{G}_{\alpha})_{\alpha \in \mathbf{A}}\) 是一族局部凸空间，\((\mathbf{F}_{\lambda})_{\lambda \in \mathbf{L}}\) 是一族向量空间。对每个 \(\lambda \in \mathbf{L}\) ，设 \(h_{\lambda}\) 是 \(\mathbf{F}_{\lambda}\) 到向量空间 E 中的线性映射；对每个 \(\lambda \in \mathbf{L}\) 及 \(\alpha \in \mathbf{J}_{\lambda}\) ，设 \(g_{\lambda \alpha}\) 是 \(\mathbf{G}_{\alpha}\) 到 \(\mathbf{F}_{\lambda}\) 中的线性映射。记 \(f_{\alpha} = h_{\lambda} \circ g_{\lambda \alpha}\) 。假设每个 \(\mathbf{F}_{\lambda}\) 带有使诸 \(g_{\lambda \alpha} (\alpha \in \mathbf{J}_{\lambda})\) 连续的最细局部凸拓扑。则 E 上使诸 \(f_{\alpha}\) 连续的最细局部凸拓扑与使诸 \(h_{\lambda}\) 连续的最细局部凸拓扑相同。

这是 S, IV, § 2.5 判别准则 CST 19 的一个特殊情形，也可以利用 prop. 5 直接证明。

局部凸终拓扑的例子。

\paragraph{I. 商空间。}

设 M 是局部凸空间 F 的一个子空间，\(\phi\) 是 F 到 F/M 上的典范映射。由于 F/M 上的商拓扑是局部凸的，并且是使 \(\phi\) 连续的所有拓扑（不论是否局部凸）中最细的一个，它也是由单独一个映射 \(\phi\) 组成的族的终局部凸拓扑。

\paragraph{II. 局部凸空间的归纳极限。}

设 A 是一个右定向有序集，\((\mathrm{E}_{\alpha}, f_{\beta\alpha})\) 是关于集 A 的一个向量空间归纳系 (A, II, § 6.2)；设 \(E = \varinjlim E_{\alpha}\)，并设 \(f_{\alpha}: E_{\alpha} \to E\) 是典范线性映射（对每个 \(\alpha \in A\) ）。假设每个 \(E_{\alpha}\) 带有一个局部凸拓扑 \(T_{\alpha}\) ，并进一步假设对 \(\alpha \leqslant \beta\) ，映射 \(f_{\beta\alpha}: E_{\alpha} \to E_{\beta}\) 连续。这时我们称族 \((\mathcal{T}_{\alpha})\) 关于线性映射族 \(f_{\alpha}\) 的终局部凸拓扑 T（相应地，带拓扑 T 的空间 E）为族 \((\mathcal{T}_{\alpha})\) 的归纳极限（相应地，系 \((\mathrm{E}_{\alpha}, f_{\beta\alpha})\) 的归纳极限空间，或径称局部凸空间 \(E_{\alpha}\) 的归纳极限空间）。回忆 E 是诸向量子空间 \(f_{\alpha}(\mathrm{E}_{\alpha})\) 的并，且当 \(\alpha \leqslant \beta\) 时，有 \(f_{\alpha}(\mathrm{E}_{\alpha}) \subset f_{\beta}(\mathrm{E}_{\beta})\) ；若给 \(f_{\alpha}(\mathrm{E}_{\alpha})\) 装备关于映射 \(f_{\alpha}\) 的终拓扑（这等同于把 \(f_{\alpha}(\mathrm{E}_{\alpha})\) 与商空间 \(E_{\alpha}/f_{\alpha}^{-1}(0)\) 等同起来），则拓扑 T 也是诸 \(f_{\alpha}(\mathrm{E}_{\alpha})\) 的拓扑所成的族关于诸典范单射的终拓扑（II, 上文 cor. 2）。此外，\(f_{\beta\alpha}\) 对 \(\alpha \leqslant \beta\) 的连续性蕴含典范单射 \(j_{\beta\alpha}: f_{\alpha}(\mathrm{E}_{\alpha}) \to f_{\beta}(\mathrm{E}_{\beta})\) 连续，因此 E 也是带有上述拓扑的诸 \(f_{\alpha}(\mathrm{E}_{\alpha})\) 关于单射 \(j_{\beta\alpha}\) 的归纳极限。

例. --- 设 X 是局部紧空间，\(E = \mathcal{K}(X; \mathbf{R})\) 是定义在 X 上、具有紧支集的有限实值连续函数所成的向量空间。对 X 的每个紧子集 K，设 \(E_{K}\) 是由那些函数 \(f \in E\) （它们满足 \(x \notin K \Rightarrow f(x) = 0\) ）组成的 E 的向量子空间。用 \(T_{K}\) 表示在 \(E_{K}\) 上诱导的拓扑，用 \(T_{u}\) 表示 X 上一致收敛的拓扑。诸拓扑 \(T_{K}\) 的归纳极限 T 比 \(T_{u}\) 细；可以证明，若 X 是仿紧的且非紧的，则 T 严格地比 \(T_{u}\) 细（参见 INT, III, 第 2 版, § 1.8）。T 的重要性在于：E 上对 T 连续的线性型恰是 X 上的实测度 (INT, III, 第 2 版, § 1.3)。

注记. --- 在最后一个例子中，\(\mathcal{T}\) 在 \(\mathrm{E}_{\mathbf{K}}\) 上诱导的拓扑与 \(\mathcal{T}_{\mathbf{K}}\) 相同，因为按定义它比 \(\mathcal{T}_{\mathbf{K}}\) 粗，而且由于 \(\mathcal{T}\) 比 \(\mathcal{T}_{u}\) 细，\(\mathcal{T}\) 在 \(\mathrm{E}_{\mathbf{K}}\) 上诱导的拓扑比 \(\mathcal{T}_{u}\) 诱导的拓扑细，也就是说比 \(\mathcal{T}_{\mathbf{K}}\) 细。

上述推理立即推广到局部凸拓扑 \((\mathcal{T}_{\alpha})\) 的归纳极限的情形，只要存在 E 上的一个局部凸拓扑 \(\mathcal{T}'\) ，使 \(\mathcal{T}_{\alpha}\) 是在 \(\mathrm{E}_{\alpha}\) 上由 \(\mathcal{T}'\) 诱导的拓扑。

更一般地，当假设 \(E_{\beta} \subset E_{\alpha}\) 且 \(T_{\beta}\) 是 \(T_{\alpha}\) 诱导的拓扑时，可以问在什么条件下 T 把 \(T_{\alpha}\) 诱导到每个 \(E_{\alpha}\) 上。一般说来并非如此 (II, 第 80 页, exerc. 26)；但我们将在以下各节中看到出现这种情况的两种重要情形。

